\section{相关工作}\label{sec:related}

\subsection{高能效语义通信}
高能效语义通信同时考虑信息提取、处理与传输的开销。
Zheng 等研究空中辅助边缘网络，针对语义传输与编码器更新提出激励机制和分布式学习方法
\cite{zheng2024aerialee}。
PSCom 在时延、功率、带宽及语义压缩约束下最小化通信与计算能耗
\cite{dai2025pscom}。
基于速率分裂的方案联合分配通信资源与语义压缩等级\cite{xu2025rsmagreen}。
相关建模还涉及计入飞行能耗的无人机网络\cite{wang2026rsma}，
以及考虑知识更新开销的可见光通信系统\cite{zhao2026vlc}。
这些研究说明，减少传输数据所带来的收益，需要与语义处理的计算开销一起评估。

对于基于 VLM 的服务，推理工作负载也会影响能耗核算。
Li 等针对机载视觉问答，在准确率约束下优化图像分辨率、传输功率与无人机轨迹
\cite{li2026onboard}。
Zhan 等的设备功耗测量表明，在其评估的 VLM 配置中，自回归输出生成是主要能耗来源之一
\cite{zhan2026seeing}。
其结果说明，仅减少视觉输入可能只能带来有限的推理节能。
这些发现促使我们考虑信息传输之后的处理开销。
本文评估的表示选择还会改变回答问题所需的计算：
检测信息对应符号解码，图像对应 VLM 推理。
由于能源供给不同，核算中区分无人机侧与边缘侧的开销。

\subsection{面向任务的视觉通信}
学习型图像传输是降低通信开销的一种方式。
深度联合信源信道编码（DJSCC）将图像直接映射为信道符号，并在接收端重建图像
\cite{bourtsoulatze2019djscc}。
基于注意力的编码适应信道条件\cite{yang2023adjscc}，
预测式自适应编码则根据估计的重建质量选择传输速率\cite{zhang2023padc}。
这些方法改善图像表示的传输。
面向任务的方案进一步依据下游任务的需求决定应传输什么。

针对无线视觉问答，MU-DeepSC 通过联合设计收发机，学习相关联的图像与问题表示
\cite{xie2022mudeepscwcl,xie2021mudeepsc}。
面向目标的视觉问答通信根据问题排序目标框与场景图三元组，并采用答案推理器
\cite{liu2024gosg}。
其针对不同问题类别的表示比较，与证据选择直接相关。
VideoQA-SC 通过带宽自适应编码，从传输的语义信息中回答视频问题，无需重建视频
\cite{guo2025videoqasc}。
这些方法已经通过最终答案的质量评价通信效果，而不要求准确的像素重建。

其他研究根据任务调整表示或其传输保护。
Park 和 Yoon 比较目标信息、布局与场景图在不同视觉任务中的表现，并过滤冗余图内容
\cite{park2025transmit}。
TONIC 根据 token 与任务的相关程度分配不等保护，
结合接收端门控与 token 补全完成图像分类\cite{liu2026tonic}。
在使用预训练多模态模型的研究中，面向查询的编码选择用于查询相关图像重建的特征
\cite{huang2026query}；
Du 等则为基于 LLaVA 的车载助手确定图像区域与图像块传输的优先级
\cite{jiang2025lmmvn}。
这些工作为依任务调整表示与传输提供了基础。
本文具体研究同一结构化视觉问答工作负载中两条完整问答路径之间的选择。
每条路径具有各自的表示、解码器、准确率与能耗。
因此，根据所选解码器的不同，较小载荷可能同时节省传输能量并避免 VLM 推理。

\subsection{自适应推理与证据选择}
自适应推理调整计算执行位置，或传送至服务器的信息量。
Neurosurgeon 在移动设备与云端之间划分神经网络的执行，以降低时延与移动端能耗
\cite{kang2017neurosurgeon}。
信息瓶颈方法为面向任务的边缘推理学习紧凑表示\cite{shao2022taskoriented}。
针对视觉问答，INAR-VL 利用图像与文本复杂度信息，在互补的边缘与云端 VLM 之间路由请求
\cite{sabanovic2026inarvl}。
这实现了根据输入选择模型与执行位置。

自适应证据传输提供了另一类相关控制变量。
Song 和 Kim 首先传输缩略图，供服务器端 VLM 推理\cite{song2026collaborative}。
服务器再根据输出不确定性，决定是否请求额外图像区域。
渐进式语义通信将视觉 token 压缩为可逐步细化的表示，
在不同通信预算下提供给现成的 VLM\cite{hsu2026progressive}。
SAGE 在硬性上行预算下选择互补视觉证据，并评估相应的分类准确率\cite{choi2026sage}。
在遥感视觉问答中，Prompt-RSVQA 将提取的视觉上下文转为文本，供语言模型使用
\cite{chappuis2022promptrsvqa}。
因此，自适应证据选择与结构化上下文都是已有的推理支持方式。

本文将发送端证据选择与问答机制选择结合起来。
发送端在传输前，选择供符号解码器使用的检测记录，或供冻结 VLM 使用的图像。
该决策不要求接收端先执行一次 VLM 推理，也不依赖额外证据请求。
我们结合已记录结果的回放、GPU 功耗测量，以及传输和机载检测开销模型，
评估其能耗与准确率之间的权衡。
这使证据级路由区别于模型部署位置的选择，
也区别于在每个问题上仍运行 VLM、仅调整其输入的方式。
